%
% Discussion.tex
%
% Copyright (C) 2025 by Inspirefly.
%
% 2.305 GHz Circular Patch Antenna Documentation
%

\chapter{Discussion}

\section{Interpretation of Results}

The analytical cavity model predicts the required radius to within 0.5~mm of the HFSS-tuned value, which confirms that Eq.~\ref{eq:fres} with the Shen fringing correction (Eq.~\ref{eq:ae}) is adequate for initial sizing at 2.305~GHz on FR-408HR. The early $\lambda_g/2$ square-patch estimate gave a correct order of magnitude but does not describe the circular $\mathrm{TM}_{11}$ resonance; its retention in git is useful only as a record of the initial bounding exercise.

The MATLAB particle-swarm run recorded in \texttt{best.txt} did not produce a viable antenna. The combination of high weights on gain and axial ratio with a single-layer, lossless substrate model drove the optimizer toward a small truncated-corner patch with a feed near the geometric centre, where the $\mathrm{TM}_{11}$ mode has a current minimum and input resistance near zero. The resulting RL${}=0$~dB and $-6.97$~dBi gain are therefore not a failure of the concept but a consequence of an under-constrained search space and an idealized feed model. This outcome motivated the shift to direct parametric sweeps in HFSS, where the feed radius $\rho_0$ and patch radius $a$ are varied independently and the full multilayer stack is represented.

HFSS results show that a single orthogonal feed pair, when driven in quadrature externally, achieves the two requirements that the MATLAB model could not satisfy simultaneously: 50~$\Omega$ match and axial ratio below 3~dB. The hexagonal patch variant confirmed that reducing symmetry introduces mode splitting; the additional $\mathrm{TM}_{21}$ resonance near 3.8~GHz coupled into the operating band and increased cross-polarization, which is why that branch was not carried into KiCAD.

\section{Substrate and Stackup Trade-offs}

Moving from Rogers RO4350B to FR-408HR increased dielectric loss from $\tan\delta \approx 0.0037$ to $0.0092$, which reduces radiation efficiency by an estimated 0.3--0.5~dB and narrows the matched bandwidth by $\sim 10\%$. The benefit is a single qualified multilayer process for the entire spacecraft and availability of 0.1~mm prepreg for controlled impedance elsewhere on the bus. The final stack places most of the dielectric thickness (0.98~mm core) between the two internal copper layers, not directly under the patch; the field is therefore concentrated in the 0.1~mm prepreg above In1.Cu, which preserves the thin-substrate approximation and limits surface-wave excitation. The total thickness 1.3414~mm remains compatible with the PocketQube rail envelope.

A thin single-layer model ($\epsilon_r=4.0$, $h=0.0994$~mm) was deliberately used in MATLAB to underestimate bandwidth, so any design that meets the RL${}<-10$~dB criterion in that model will have margin in the fabricated stack. The measured $\epsilon_r$ of FR-408HR varies $\pm 0.05$ across lots; a $\pm 0.05$ change shifts $f_0$ by $\mp 18$~MHz ($\sim 0.8\%$), which is within the 20~MHz target window but requires a final HFSS retune if a different lot is used. No post-fabrication tuning stubs were added; frequency trimming, if needed, will be done by a single HFSS re-spin of radius $a$ rather than by matching-network components.

\section{Feed Strategy and the External Divider}

The decision to use two coax feeds with an external divider (commit \texttt{5554f5c}: ``Use external power divider: APS-02-WBM-LP-DF-CC'') reflects three constraints. First, an on-board branch-line hybrid at 2.305~GHz would occupy $\approx 20 \times 20$~mm, consuming area needed for the patch ground plane. Second, the external divider allows the same board to be tested in linear polarization (one port terminated) and circular polarization (quadrature drive) without layout changes. Third, the Molex 732511350 U.FL-compatible connector is qualified for vibration and has a smaller keep-out than an SMA, which preserves ground-plane continuity near the patch edge.

The trade is an additional 0.4~dB insertion loss and the need for phase-matched cables. Cable phase error of $\pm 5^{\circ}$ degrades axial ratio by $\approx 0.4$~dB, which is acceptable against the 3~dB requirement. The SMA edge-mount footprint is retained in the schematic as a laboratory alternative where cable loss can be de-embedded with a VNA.

\section{Version Control and Reproducibility}

Git history shows a deliberate narrowing from general-purpose explorations (monopole, hexagonal) to the circular patch V3 that is now in KiCAD. Each geometry change is tagged by an AEDT file commit, so the HFSS source for any KiCAD backup ZIP can be traced: \texttt{CircularPatchV3.aedt} at \texttt{235c732} is the parent of \texttt{KiCAD/CircularPatchV3.dxf} and therefore of every \texttt{antennaBoard-2026-0x\_xxxxxx.zip} after 2026-03-01. The DXF-to-KiCAD Python converter is pinned in the repository, so the polygon conversion is reproducible to $10~\mu$m. The \texttt{LOCK.txt} convention prevented concurrent KiCAD edits from producing merge conflicts in the binary \texttt{.kicad\_pcb} file, which is not diffable.

Remaining uncertainty lies in fabrication tolerances: etch undercut of $\pm 0.05$~mm on a 19~mm radius shifts $f_0$ by $\pm 6$~MHz, and via drill wander of $\pm 0.05$~mm shifts the feed resistance by $\pm 3$~$\Omega$. Both are within the matching network margin if $S_{11} < -14$~dB is achieved in HFSS, but they argue for a first-article VNA measurement and, if needed, a single corrective spin of $a$ by 0.1~mm increments.

\section{Limitations}

No anechoic gain or axial-ratio measurements are reported in this revision; all pattern data are from HFSS far-field calculations. The MATLAB model does not include the multilayer stack or connector parasitics. The KiCAD board has not yet been fabricated at the final 1.3414~mm stack, so return loss and inter-port isolation are predicted only. These gaps define the verification work in Chapter~5.
